% TEMPLATE_PLACEHOLDER: this is a formatting guide, not a scientific article.
\section{Постановка задачи}
Здесь следует определить исследовательскую задачу, область применимости и
необходимые исходные условия. Статья должна быть понятна читателю без обращения
к отдельной учебной или научно-исследовательской работе. Включите определения
объектов и обозначений, которые используются далее. Укажите конкретное
ограничение существующих подходов и объясните, какой результат требуется получить.

Этот проект содержит пример оформления, который нужно заменить текстом статьи.
Образец ссылки на литературу: \cite{sample}. Источник взят из официального
шаблона; замените его литературой, использованной в вашем исследовании.

\section{Предлагаемый подход}
В этой части опишите метод, модель или алгоритм решения поставленной задачи.
Введите необходимые величины, их смысл и допустимые значения. Сформулируйте
исходные предположения и ограничения. Отделяйте предложенные решения от сведений,
заимствованных из литературы. Ссылку размещайте рядом с утверждением, которое
она подтверждает. Не добавляйте технические подробности, не влияющие на понимание
метода или оценку его применимости.

В примере формулы $x_i$ обозначает вклад элемента, $n$ --- число элементов.
Единственную формулу без ссылки можно оформить без номера.
\begin{equation}
  S = \sum_{i=1}^{n} x_i .
  \label{eq:sum}
\end{equation}
Ссылка на формулу~\eqref{eq:sum} должна использовать автоматически назначенный
номер. Не вписывайте номера формул, рисунков или таблиц вручную. Если объектов
одного вида несколько, они должны быть пронумерованы последовательно.

\section{Обоснование и пример}
Приведите обоснование предложенного подхода или результаты эксперимента.
Если используется расчётный пример, явно определите исходные данные и
покажите, как из них получается результат. Демонстрационные данные обозначайте
как демонстрационные. Не выдавайте предполагаемый эффект за измеренный.
Сопоставление с исходным подходом должно проводиться в одинаковых условиях.

Название таблицы располагается сверху справа, как в этом образце.
При единственной таблице номер скрывается автоматически после повторной сборки.
\begin{table}[htbp]
  \caption{Пример структуры сопоставления}
  \label{tab:example}
  \centering
  \begin{tabular*}{\linewidth}{@{\extracolsep{\fill}}ll@{}}
    \toprule
    Элемент анализа & Содержание \\
    \midrule
    Исходные условия & Описание предположений \\
    Предлагаемое решение & Описание метода \\
    Проверка & Обоснование свойств \\
    \bottomrule
  \end{tabular*}
\end{table}

Рисунки должны быть чёрно-белыми, подписи располагаются снизу по центру.
Проверьте читаемость при печати в А5. Примеры вставки объектов находятся
в файле \textnormal{examples/snippets.tex}.

\section{Заключение}
Сформулируйте конкретный результат, соответствующий поставленной задаче.
Укажите существенные ограничения и условия применимости. Выводы должны следовать
из приведённых рассуждений или данных. Объём готовых материалов составляет ровно
две полные страницы А5, включая аннотацию, рисунки, таблицы и список литературы.
Окончательную заполненность второй страницы проверяйте визуально: пустая страница
или несколько строк на ней не удовлетворяют требованию полноты.
